\subsection{Интерпретация}

\textbf{Что повышает зарплату.} Самые сильные прибавки дают навыки интеграционной и архитектурной работы, добавленные к базовому стеку системного аналитика:
\begin{itemize}
\item \textbf{Kafka}: $+139$~тыс. к BPMN, $+133$~тыс. к SQL, $+108$~тыс. к UML; в тройках $+100\ldots105$~тыс. к парам BPMN+UML, BPMN+SQL, SQL+UML. Медиана вакансий с Kafka --- около 300~тыс. на руки.
\item \textbf{Микросервисная архитектура} ($+120{,}5$~тыс. к SQL) и \textbf{системная интеграция} ($+100$~тыс. к SQL), \textbf{XML/JSON} к базам данных и постановке задач ($+90\ldots100$~тыс.).
\item \textbf{A/B-тесты} к Python ($+105$~тыс.) и к паре Python+SQL ($+100$~тыс.) --- продуктовая аналитика с экспериментами.
\item \textbf{ERP-системы на базе 1С} и \textbf{бухгалтерский учёт} к «сбору и анализу информации» ($+120$~тыс.): узкая и дорогая ниша внедрения ERP.
\end{itemize}

\textbf{Что понижает.} Сильнее всего медиану тянут вниз \emph{общие} формулировки, добавленные к технической базе: «работа с базами данных» к REST ($-120$~тыс.), «аналитический склад ума» к системному анализу ($-95$~тыс.), «оптимизация бизнес-процессов» к UML ($-89$~тыс.), MS Excel, анализ рынка, аналитические исследования ($-69\ldots80$~тыс.). В тройках --- «информационные технологии», «аналитическое мышление», «техническая документация» к дорогим парам с REST/SOAP ($-90\ldots106$~тыс.). Неожиданно понижает и \textbf{Python} в тройке с BPMN+SQL ($-95$~тыс.): такой набор характерен не для системных аналитиков, а для начинающих аналитиков данных и бизнес-аналитиков.

\begin{callout}
\textbf{Главное объяснение --- опыт, а не навык сам по себе.} В последнем столбце таблиц видно: у повышающих пар доля вакансий с опытом от 3 лет 43--97\% (у большинства выше 55\%) против 21--52\% у базы, а у понижающих --- 3--50\% при 34--62\% у базы. Технические интеграционные навыки работодатели перечисляют в вакансиях для опытных системных аналитиков, а общие и «мягкие» формулировки --- в вакансиях для начинающих и в массовых объявлениях (например, стажёрские вакансии работодателя 6093775 с медианой 65\,250~₽). Поэтому прибавка --- это описание рынка, а не обещание: выучив Kafka без опыта, зарплату в 300~тыс. не получить.
\end{callout}
